        %%******************************************%%
        %%                                          %%
        %%	           Tesi di laurea               %%
        %%                                          %%
        %%             14 Luglio 2017               %%
        %%                                          %%
        %%******************************************%%


% I seguenti commenti speciali impostano:
% 1. 
% 2. PDFLaTeX come motore di composizione;
% 3. tesi.tex come documento principale;
% 4. il controllo ortografico italiano per l'editor.

% !TEX encoding = UTF-8
% !TEX TS-program = pdflatex
% !TEX root = tesi.tex
% !TEX spellcheck = it-IT

\documentclass[11pt,                    % corpo del font principale
               a4paper,                 % carta A4
               twoside,                 % impagina per fronte-retro
               openright,               % inizio capitoli a destra
               english,                 
               italian,                 
               ]{book}    

\usepackage[utf8]{inputenc}             % codifica di input; anche [latin1] va bene
                                        % NOTA BENE! va accordata con le preferenze dell'editor

%**************************************************************
% Importazione package
%************************************************************** 

%\usepackage{amsmath,amssymb,amsthm}    % matematica

\usepackage[english, italian]{babel}    % per scrivere in italiano e in inglese;
                                        % l'ultima lingua (l'italiano) risulta quella predefinita

\usepackage{bookmark}                   % segnalibri

\usepackage{titlesec} 

\usepackage{caption}                    % didascalie

\usepackage{chngpage,calc}              % centra il frontespizio

\usepackage{csquotes}                   % gestisce automaticamente i caratteri (")

\usepackage{emptypage}                  % pagine vuote senza testatina e piede di pagina

\usepackage{epigraph}					          % per epigrafi

\usepackage{eurosym}                    % simbolo dell'euro

\usepackage[T1]{fontenc}                % codifica dei font:
                                        % NOTA BENE! richiede una distribuzione *completa* di LaTeX

%\usepackage{indentfirst}               % rientra il primo paragrafo di ogni sezione

\usepackage{graphicx}                   % immagini

\usepackage{enumerate}                  % liste numerate

\usepackage{hyperref}                   % collegamenti ipertestuali



\usepackage[binding=5mm]{layaureo}      % margini ottimizzati per l'A4; rilegatura di 5 mm

\usepackage{listings}                   % codici

\usepackage{microtype}                  % microtipografia

\usepackage{mparhack,relsize}  % finezze tipografiche

\usepackage{nameref}                    % visualizza nome dei riferimenti                                      

\usepackage[font=small]{quoting}        % citazioni

\usepackage{subfig}                     % sottofigure, sottotabelle

\usepackage[italian]{varioref}          % riferimenti completi della pagina

\usepackage[dvipsnames]{xcolor}         % colori

\usepackage{booktabs}                   % tabelle                                       
\usepackage{tabularx}                   % tabelle di larghezza prefissata                                    
\usepackage{longtable}                  % tabelle su più pagine                                        
\usepackage{ltxtable}                   % tabelle su più pagine e adattabili in larghezza

\usepackage[automake, toc, acronym]{glossaries}   % glossario
                                        % per includerlo nel documento bisogna:
                                        % 1. compilare una prima volta tesi.tex;
                                        % 2. eseguire: makeindex -s tesi.ist -t tesi.glg -o tesi.gls tesi.glo
                                        % 3. eseguire: makeindex -s tesi.ist -t tesi.alg -o tesi.acr tesi.acn
                                        % 4. compilare due volte tesi.tex.

\usepackage[backend=biber,style=verbose-ibid,hyperref,backref]{biblatex}
                                        % eccellente pacchetto per la bibliografia; 
                                        % produce uno stile di citazione autore-anno; 
                                        % lo stile "numeric-comp" produce riferimenti numerici
                                        % per includerlo nel documento bisogna:
                                        % 1. compilare una prima volta tesi.tex;
                                        % 2. eseguire: biber tesi
                                        % 3. compilare ancora tesi.tex.

%**************************************************************
% file contenente le impostazioni della tesi
%**************************************************************

%**************************************************************
% Frontespizio
%**************************************************************
\newcommand{\myName}{Mario Rossi}                                    % autore
\newcommand{\myTitle}{Titolo della mia tesi di laurea}                    
\newcommand{\myDegree}{Tesi di laurea triennale}                % tipo di tesi
\newcommand{\myUni}{Università degli Studi di Padova}           % università
\newcommand{\myFaculty}{Corso di Laurea in Qualcosa}         % facoltà
\newcommand{\myDepartment}{Dipartimento di Matematica "Tullio Levi-Civita"}          % dipartimento
\newcommand{\myProf}{NomeCognome relatore}                                % relatore
\newcommand{\myLocation}{Padova}                                % dove
\newcommand{\myAA}{2016-2017}                                   % anno accademico
\newcommand{\myTime}{Luglio 2017}                                  % quando


%**************************************************************
% Impostazioni di impaginazione
% see: http://wwwcdf.pd.infn.it/AppuntiLinux/a2547.htm
%**************************************************************

\setlength{\parindent}{14pt}   % larghezza rientro della prima riga
\setlength{\parskip}{0pt}   % distanza tra i paragrafi


%**************************************************************
% Impostazioni dello stile per i chapter
%**************************************************************

\titleformat{\chapter}[display]
	{\bfseries\Large}
	{\filright \MakeUppercase{\chaptertitlename} \Huge\thechapter}
	{1ex}
	{\titlerule\vspace{1ex}\filright}
	[\vspace{1ex}\titlerule]

%**************************************************************
% Impostazioni di biblatex
%**************************************************************
\bibliography{bibliografia} % database di biblatex 

\defbibheading{bibliography}
{
    \cleardoublepage
    \phantomsection 
    \addcontentsline{toc}{chapter}{\bibname}
    \chapter*{\bibname\markboth{\bibname}{\bibname}}
}

\setlength\bibitemsep{1.5\itemsep} % spazio tra entry

\DeclareBibliographyCategory{opere}
\DeclareBibliographyCategory{web}

\addtocategory{opere}{womak:lean-thinking}
\addtocategory{web}{site:agile-manifesto}

\defbibheading{opere}{\section*{Riferimenti bibliografici}}
\defbibheading{web}{\section*{Siti Web consultati}}


%**************************************************************
% Impostazioni di caption
%**************************************************************
\captionsetup{
    tableposition=top,
    figureposition=bottom,
    font=small,
    format=hang,
    labelfont=bf
}

%**************************************************************
% Impostazioni di glossaries
%**************************************************************

%**************************************************************
% Glossario
%**************************************************************
\renewcommand{\glossaryname}{Glossario}

\newglossaryentry{apig}
{
    name=\glslink{apig}{API},
    text=API,
    sort=api,
    description={In informatica con il termine \emph{Application Programming Interface API} (ing. interfaccia di programmazione di un'applicazione) si indica ogni insieme di procedure disponibili al programmatore, di solito raggruppate a formare un set di strumenti specifici per l'espletamento di un determinato compito all'interno di un certo programma. La finalità è ottenere un'astrazione, di solito tra l'hardware e il programmatore o tra software a basso e quello ad alto livello semplificando così il lavoro di programmazione}
}

% \newglossaryentry{umlg}
% {
%     name=\glslink{uml}{UML},
%     text=UML,
%     sort=uml,
%     description={in ingegneria del software \emph{UML, Unified Modeling Language} (ing. linguaggio di modellazione unificato) è un linguaggio di modellazione e specifica basato sul paradigma object-oriented. L'\emph{UML} svolge un'importantissima funzione di ``lingua franca'' nella comunità della progettazione e programmazione a oggetti. Gran parte della letteratura di settore usa tale linguaggio per descrivere soluzioni analitiche e progettuali in modo sintetico e comprensibile a un vasto pubblico}
% }

\newglossaryentry{teamviewerg}
{
    name=\glslink{teamviewerg}{Team Viewer},
    text=team Viewer,
    sort=teamviewer,
    description={software gratuito per il controllo remoto dei computer. Si tratta, dunque, di un programma che consente di dirigere i computer a distanza e di ricevere assistenza remota tramite Internet}
}

\newglossaryentry{systemintegratorg}
{
    name=\glslink{systemintegratorg}{System integrator},
    text=System integrator,
    sort=systemintegrator,
    description={Con il termine inglese system integrator viene indicata una azienda (o uno specialista) che si occupa dell'integrazione di sistemi.\\ Il compito del system integrator è quello di far dialogare impianti diversi tra di loro allo scopo di creare una nuova struttura funzionale che possa utilizzare sinergicamente le potenzialità degli impianti d'origine e creando quindi funzionalità originariamente non presenti}
}

\newglossaryentry{freemiumg}
{
    name=\glslink{freemiumg}{Freemium},
    text=fremium,
    sort=Freemium,
    description={Modello economico consistente nel proporre al consumatore due versioni dello stesso prodotto, una più semplice, gratuita, ed una più ricca, a pagamento}
}

\newglossaryentry{filehostingg}
{
    name=\glslink{filehostingg}{File hosting},
    text=file hosting,
    sort=File hosting,
    description={Un servizio di \emph{file hosting} (in italiano: archiviazione di file), chiamato anche \emph{cyberlocker}, è un servizio di archiviazione su Internet appositamente progettato per ospitare i file degli utenti, permettendo loro di caricare file che possono poi essere scaricati da altri utenti}
}

\newglossaryentry{storageg}
{
    name=\glslink{storageg}{Storage Cloud},
    text=storage Cloud,
    sort=Storage Cloud,
    description={Il Cloud Storage è un modello di conservazione dati su computer in rete dove i dati stessi sono memorizzati su molteplici server virtuali generalmente ospitati presso strutture di terze parti o su server dedicati}
}

\newglossaryentry{filesharingg}
{
    name=\glslink{filesharingg}{File sharing},
    text=file sharing,
    sort=File sharing,
    description={In informatica il file sharing è la condivisione di file all'interno di una rete di calcolatori e tipicamente utilizza architetture client-server o peer-to-peer. Queste reti possono permettere di ricercare un file in particolare per mezzo di un URI (Universal Resource Identifier), di individuare più copie dello stesso file nella rete per mezzo di impronte (hash), di eseguire lo scaricamento da più fonti contemporaneamente, di riprendere lo scaricamento del file dopo un'interruzione.\\
    I programmi di filesharing sono utilizzati direttamente o indirettamente per trasferire file da un computer ad un altro su Internet, o su reti aziendali Intranet. Questa condivisione ha dato origine al modello peer-to-peer}
}

\newglossaryentry{helpdeskg}
{
    name=\glslink{helpdeskg}{Help desk},
    text=help desk,
    sort=Help desk,
    description={L'help desk è una risorsa destinata a fornire supporto, all'utente o al cliente, relativamente a prodotti o servizi informatici ed elettronici. Lo scopo dell'Help desk è di risolvere problemi o fornire indicazioni su prodotti come computer, apparecchiature elettroniche o software}
}

\newglossaryentry{embedg}
{
    name=\glslink{embedg}{Embed widget},
    text=embed widget,
    sort=Embed widget,
    description={Con un widget per l'embedding è possibile, tramite un link. aggiungere alle vostre pagine elementi multimediali, come video o brani musicali, provenienti da altri siti Internet}
}

\newglossaryentry{paasg}
{
    name=\glslink{paasg}{Platform-as-a-Service},
    text=Platform-as-a-Service,
    sort=Platform-as-a-Service,
    description={\emph{"Platform as a Service"}, spesso semplicemente chiamata \emph{PaaS}, è una categoria di cloud computing che fornisce agli sviluppatori una piattaforma e un ambiente per costruire applicazioni e servizi su Internet}
}

\newglossaryentry{frameworkg}
{
    name=\glslink{frameworkg}{Framework},
    text=framework,
    sort=Framework,
    description={in informatica e specificatamente nello sviluppo software, è un'architettura logica di supporto (spesso un'implementazione logica di un particolare design pattern) su cui un software può essere progettato e realizzato, facilitandone lo sviluppo da parte del programmatore}
}

\newglossaryentry{rsag}
{
    name=\glslink{rsag}{RSA},
    text=RSA,
    sort=RSA,
    description={In crittografia la sigla RSA indica un algoritmo di crittografia asimmetrica, inventato nel 1977 da Ronald Rivest, Adi Shamir e Leonard Adleman utilizzabile per cifrare o firmare informazioni}
}

\newglossaryentry{sdkg}
{
    name=\glslink{sdkg}{SDK},
    text=SDK,
    sort=SDK,
    description={Acronimo di \emph{software development kit} (SDK, traducibile in italiano come "pacchetto di sviluppo per applicazioni"), in informatica, indica genericamente un insieme di strumenti per lo sviluppo e la documentazione di software}
}

\newglossaryentry{commitg}
{
    name=\glslink{commitg}{Commit},
    text=commit,
    sort=Commit,
    description={Nei sistemi di controllo di versione, un commit è l'operazione atomica che aggiunge le modifiche più recenti a (parte del) codice sorgente nel repository, rendendo queste modifiche parte della revisione dell'head del repository}
}

\newglossaryentry{buildautomationg}
{
    name=\glslink{buildautomationg}{Build automation},
    text=build automation,
    sort=Build automation,
    description={L'Automazione dello sviluppo (in inglese: \emph{build automation}), in informatica, è l'atto di scrivere o automatizzare un'ampia varietà di compiti che gli sviluppatori software fanno nelle loro attività quotidiane di sviluppo come: compilazione del codice sorgente in codice binario, pacchettizzazione del codice binario, esecuzione di test, deployment di sistemi di produzione, creazione di documentazione e/o note di rilascio}
}

\newglossaryentry{bottomupg}
{
    name=\glslink{bottomupg}{Bottom-up},
    text=bottom-up,
    sort=Bottom-up,
    description={Nella progettazione bottom-up, le parti individuali del sistema sono specificate in dettaglio, e poi connesse tra loro in modo da formare componenti più grandi, a loro volta interconnesse fino a realizzare un sistema completo. Le strategie basate sul flusso informativo bottom-up sembrano potenzialmente necessarie e sufficienti, poiché basate sulla conoscenza di tutte le variabili in grado di condizionare gli elementi del sistema}
}

\newglossaryentry{ciclog}
{
    name=\glslink{ciclog}{Complessità ciclomatica},
    text=complessità ciclomatica,
    sort=Complessità ciclomatica,
    description={La Complessità Ciclomatica (o complessità condizionale) è una metrica software. Sviluppata da Thomas J. McCabe nel 1976, è utilizzata per misurare la complessità di un programma. Misura direttamente il numero di cammini linearmente indipendenti attraverso il grafo di controllo di flusso.\\
    I nodi del grafo corrispondono a gruppi indivisibili di istruzioni, mentre gli archi connettono due nodi se il secondo gruppo di istruzioni può essere eseguito immediatamente dopo il primo gruppo. La complessità ciclomatica può, inoltre, essere applicata a singole funzioni, moduli, metodi o classi di un programma}
}

\newglossaryentry{clientserverg}
{
    name=\glslink{clientserverg}{Client-server},
    text=client-server,
    sort=Client-server,
    description={In informatica il termine sistema \emph{client-server} (letteralmente cliente-serviente) indica un'architettura di rete nella quale genericamente un computer client o terminale si connette ad un server per la fruizione di un certo servizio, quale ad esempio la condivisione di una certa risorsa hardware/software con altri client, appoggiandosi alla sottostante architettura protocollare}
}
 % database di termini
\makeglossaries


%**************************************************************
% Impostazioni di graphicx
%**************************************************************
\graphicspath{{immagini/}} % cartella dove sono riposte le immagini


%**************************************************************
% Impostazioni di hyperref
%**************************************************************
\hypersetup{
    %hyperfootnotes=false,
    %pdfpagelabels,
    %draft,	% = elimina tutti i link (utile per stampe in bianco e nero)
    colorlinks=true,
    linktocpage=true,
    pdfstartpage=1,
    pdfstartview=FitV,
    % decommenta la riga seguente per avere link in nero (per esempio per la stampa in bianco e nero)
    %colorlinks=false, linktocpage=false, pdfborder={0 0 0}, pdfstartpage=1, pdfstartview=FitV,
    breaklinks=true,
    pdfpagemode=UseNone,
    pageanchor=true,
    pdfpagemode=UseOutlines,
    plainpages=false,
    bookmarksnumbered,
    bookmarksopen=true,
    bookmarksopenlevel=1,
    hypertexnames=true,
    pdfhighlight=/O,
    %nesting=true,
    %frenchlinks,
    urlcolor=webbrown,
    linkcolor=RoyalBlue,
    citecolor=webgreen,
    %pagecolor=RoyalBlue,
    %urlcolor=Black, linkcolor=Black, citecolor=Black, %pagecolor=Black,
    pdftitle={\myTitle},
    pdfauthor={\textcopyright\ \myName, \myUni, \myFaculty},
    pdfsubject={},
    pdfkeywords={},
    pdfcreator={pdfLaTeX},
    pdfproducer={LaTeX}
}

%**************************************************************
% Impostazioni di itemize
%**************************************************************

%\renewcommand{\labelitemi}{$\bullet$}
%\renewcommand{\labelitemii}{$\cdot$}
%\renewcommand{\labelitemiii}{$\diamond$}
%\renewcommand{\labelitemiv}{$\ast$}


%**************************************************************
% Impostazioni di listings
%**************************************************************
\lstset{
    language=[LaTeX]Tex,%C++,
    keywordstyle=\color{RoyalBlue}, %\bfseries,
    basicstyle=\small\ttfamily,
    %identifierstyle=\color{NavyBlue},
    commentstyle=\color{Green}\ttfamily,
    stringstyle=\rmfamily,
    numbers=none, %left,%
    numberstyle=\scriptsize, %\tiny
    stepnumber=5,
    numbersep=8pt,
    showstringspaces=false,
    breaklines=true,
    frameround=ftff,
    frame=single
} 


%**************************************************************
% Impostazioni di xcolor
%**************************************************************
\definecolor{webgreen}{rgb}{0,.5,0}
\definecolor{webbrown}{rgb}{.6,0,0}


%**************************************************************
% Altro
%**************************************************************

\newcommand{\omissis}{[\dots\negthinspace]} % produce [...]

% eccezioni all'algoritmo di sillabazione
\hyphenation
{
    ma-cro-istru-zio-ne
}

\newcommand{\sectionname}{sezione}
\addto\captionsitalian{\renewcommand{\figurename}{figura}
                       \renewcommand{\tablename}{tabella}}

\newcommand{\glsfirstoccur}{\ap{{[g]}}}

\newcommand{\intro}[1]{\emph{\textsf{#1}}} % piccola intro sotto il titolo del capitolo

%**************************************************************
% Environment per ``rischi''
%**************************************************************
\newcounter{riskcounter}                % define a counter
\setcounter{riskcounter}{0}             % set the counter to some initial value

%%%% Parameters
% #1: Title
\newenvironment{risk}[1]{
    \refstepcounter{riskcounter}        % increment counter
    \par \noindent                      % start new paragraph
    \textbf{\arabic{riskcounter}. #1}   % display the title before the 
                                        % content of the environment is displayed 
}{
    \par\medskip
}

\newcommand{\riskname}{Rischio}

\newcommand{\riskdescription}[1]{\textbf{\\Descrizione:} #1.}

\newcommand{\risksolution}[1]{\textbf{\\Soluzione:} #1.}

%**************************************************************
% Environment per ``use case''
%**************************************************************
\newcounter{usecasecounter}             % define a counter
\setcounter{usecasecounter}{0}          % set the counter to some initial value

%%%% Parameters
% #1: ID
% #2: Nome
\newenvironment{usecase}[2]{
    \renewcommand{\theusecasecounter}{\usecasename #1}  % this is where the display of 
                                                        % the counter is overwritten/modified
    \refstepcounter{usecasecounter}             % increment counter
    \vspace{10pt}
    \par \noindent                              % start new paragraph
    {\large \textbf{\usecasename #1: #2}}       % display the title before the 
                                                % content of the environment is displayed 
    \medskip
}{
    \medskip
}

\newcommand{\usecasename}{UC}

\newcommand{\usecaseactors}[1]{\textbf{\\Attori Principali:} #1. \vspace{4pt}}
\newcommand{\usecasepre}[1]{\textbf{\\Precondizioni:} #1. \vspace{4pt}}
\newcommand{\usecasedesc}[1]{\textbf{\\Descrizione:} #1. \vspace{4pt}}
\newcommand{\usecasepost}[1]{\textbf{\\Postcondizioni:} #1. \vspace{4pt}}
\newcommand{\usecasealt}[1]{\textbf{\\Scenario Alternativo:} #1. \vspace{4pt}}

%**************************************************************
% Environment per ``namespace description''
%**************************************************************

\newenvironment{namespacedesc}{
    \vspace{10pt}
    \par \noindent                              % start new paragraph
    \begin{description} 
}{
    \end{description}
    \medskip
}

\newcommand{\classdesc}[2]{\item[\textbf{#1:}] #2}                     % file con le impostazioni personali


\begin{document}
%**************************************************************
% Materiale iniziale
%**************************************************************
\frontmatter
% !TEX encoding = UTF-8
% !TEX TS-program = pdflatex
% !TEX root = ../tesi.tex
% !TEX spellcheck = it-IT

%**************************************************************
% Frontespizio 
%**************************************************************
\begin{titlepage}

\begin{center}

\begin{LARGE}
\textbf{\myUni}\\
\end{LARGE}

\vspace{10pt}

\begin{Large}
\textsc{\myDepartment}\\
\end{Large}

\vspace{10pt}

\begin{large}
\textsc{\myFaculty}\\
\end{large}

\vspace{30pt}
\begin{figure}[htbp]
\begin{center}
\includegraphics[height=6cm]{logo-unipd}
\end{center}
\end{figure}
\vspace{30pt} 

\begin{LARGE}
\begin{center}
\textbf{\myTitle}\\
\end{center}
\end{LARGE}

\vspace{10pt} 

\begin{large}
\textsl{\myDegree}\\
\end{large}

\vspace{40pt} 

\begin{large}
\begin{flushleft}
\textit{Relatore}\\ 
\vspace{5pt} 
Prof. \myProf
\end{flushleft}

\vspace{0pt} 

\begin{flushright}
\textit{Laureando}\\ 
\vspace{5pt} 
\myName
\end{flushright}
\end{large}

\vspace{40pt}

\line(1, 0){338} \\
\begin{normalsize}
\textsc{Anno Accademico \myAA}
\end{normalsize}

\end{center}
\end{titlepage} 
% !TEX encoding = UTF-8
% !TEX TS-program = pdflatex
% !TEX root = ../tesi.tex
% !TEX spellcheck = it-IT

%**************************************************************
% Colophon
%**************************************************************
\clearpage
\phantomsection
\thispagestyle{empty}

\hfill

\vfill

\noindent\myName: \textit{\myTitle,}
\myDegree,
\textcopyright\ \myTime.
% !TEX encoding = UTF-8
% !TEX TS-program = pdflatex
% !TEX root = ../tesi.tex
% !TEX spellcheck = it-IT

%**************************************************************
% Dedica
%**************************************************************
\cleardoublepage
\phantomsection
\thispagestyle{empty}
\pdfbookmark{Dedica}{Dedica}

\vspace*{3cm}

\begin{center}
	\emph{Qui andrà scritta la nostra dolcissima dedica}
\end{center}

\medskip

\noindent\textit{}
\hfill \emph{Vostro nome}


% !TEX encoding = UTF-8
% !TEX TS-program = pdflatex
% !TEX root = ../tesi.tex
% !TEX spellcheck = it-IT

%**************************************************************
% Sommario
%**************************************************************
\cleardoublepage
\phantomsection
\pdfbookmark{Sommario}{Sommario}
\begingroup
\let\clearpage\relax
\let\cleardoublepage\relax
\let\cleardoublepage\relax

\chapter*{Sommario}

Il presente documento descrive il lavoro svolto durante il periodo di stage, della durata di circa trecento ore, dal laureando Nome Cognome presso l'azienda NomeAzienda S.p.A.
Gli obbiettivi da raggiungere erano molteplici.\\
In primo luogo era richiesta l'analisi delle peculiarità e delle funzonalità di \emph{tecnologia/e} al fine di comprendere meglio questa piattaforma.
In secondo luogo era richiesta l'implementazione di una autorizzazione lato back-end server-to-server.
Questo tipo di autenticazione . . .
Terzo ed ultimo obbiettivo era la realizzazione di alcuni servizi, sempre lato back-end, per una loro prevedibile integrazione con alcuni processi documentali aziendali.

%\vfill
%
%\selectlanguage{english}
%\pdfbookmark{Abstract}{Abstract}
%\chapter*{Abstract}
%
%\selectlanguage{italian}

\endgroup			

\vfill


% !TEX encoding = UTF-8
% !TEX TS-program = pdflatex
% !TEX root = ../tesi.tex
% !TEX spellcheck = it-IT

%**************************************************************
% Ringraziamenti
%**************************************************************
\cleardoublepage
\phantomsection
\pdfbookmark{Ringraziamenti}{ringraziamenti}

\begin{flushright}{
	\slshape    
	''Citazione famosa o frase che ci rappresenta''} \\ 
	\medskip
    --- Autore
\end{flushright}


\bigskip

\begingroup
\let\clearpage\relax
\let\cleardoublepage\relax
\let\cleardoublepage\relax

\chapter*{Ringraziamenti}

\noindent \textit{Innanzitutto, vorrei esprimere la mia gratitudine al Prof. Nome Cognome, relatore della mia tesi, per l'aiuto e il sostegno fornitomi durante la stesura del lavoro.}\\

\noindent \textit{Ringrazio inoltre il mio tutor aziendale, Nome cognome tutor, per la disponibilità sempre avuta con me e l'azienda Nome per la possibilità di cimentarmi in questo progetto.}\\

\noindent \textit{Desidero ringraziare con particolare affetto i miei genitori fratelli/sorelle per il sostegno, il grande aiuto e per essermi stati vicini in ogni momento durante gli anni di studio; grazie di aver creduto in me.}\\

\noindent \textit{Ho desiderio di ringraziare poi tutti i miei più cari amici per tutti i bellissimi anni passati insieme e le mille avventure vissute!}\\

\noindent \textit{Voglio ringraziare tutti i componenti del gruppo mio Team, con i quali ho condiviso un anno di grande impegno e difficoltà e che, di sicuro, non dimenticherò mai.}\\

\noindent \textit{Un ringraziamento particolare va alla mia ragazza/ragazzo (se esiste) }\\

\noindent \textit{Voglio infine ringraziare anche me stesso, per qual si voglia motivo esistente}\\
\bigskip

\noindent\textit{\myLocation, \myTime}
\hfill \myName

\endgroup


% !TEX encoding = UTF-8
% !TEX TS-program = pdflatex
% !TEX root = ../tesi.tex
% !TEX spellcheck = it-IT

%**************************************************************
% Indici
%**************************************************************
\clearpage
\pdfbookmark{\contentsname}{tableofcontents}
\setcounter{tocdepth}{2} % profondità indice desiderata
\tableofcontents
%\markboth{\contentsname}{\contentsname} 
\clearpage

\begingroup 
    \let\clearpage\relax
    \let\cleardoublepage\relax
    \let\cleardoublepage\relax
    %*******************************************************
    % Elenco delle figure
    %*******************************************************    
    \phantomsection
    \pdfbookmark{\listfigurename}{lof}
    \listoffigures

    \vspace*{8ex}

    %*******************************************************
    % Elenco delle tabelle
    %*******************************************************
    \phantomsection
    \pdfbookmark{\listtablename}{lot}
    \listoftables
        
    \vspace*{8ex}
\endgroup

\cleardoublepage

\cleardoublepage

%**************************************************************
% Materiale principale
%**************************************************************
\mainmatter
% !TEX encoding = UTF-8
% !TEX TS-program = pdflatex
% !TEX root = ../tesi.tex
% !TEX spellcheck = it-IT

%**************************************************************
\chapter{Il contesto aziendale}
\label{cap:contesto}
%**************************************************************

\intro{Questo capitolo introduce al contesto aziendale nel quale si è svolto lo stage, i prodotti ed i servizi che l'azienda offre e l'ottica nel quale il progetto si inserisce, in relazione ad obiettivi e benefici aziendali.}

% Introduzione al contesto applicativo.\\

% \noindent Esempio di utilizzo di un termine nel glossario \\
% \gls{api}. \\

% \noindent Esempio di citazione in linea \\
% \cite{site:agile-manifesto}. \\

% \noindent Esempio di citazione nel pie' di pagina \\
% citazione\footcite{womak:lean-thinking} \\

%**************************************************************
\section{L'azienda}
\label{sec:azienda}

Parliamo dell'azienda.\\
\textbf{Azienda S.p.A.} è un'azienda che nasce nel 19mila come \gls{systemintegratorg}\glsfirstoccur.\\


\begin{figure}[!h]
	\centering
	\includegraphics[width=\textwidth]{tesi-template/immagini/cap2/logo_google_esempio.png}
	\caption{Logo azienda}
	\label{fig:google-logo}
\end{figure}

\newpage

Obbiettivi:
\begin{itemize}
		\item \underline{\textit{De1}}\label{obj:de1}: Progettazione e realizzazione di servizi che permettano la manipolazione in locale dei commenti associati a documenti remoti su \emph{Box}, in particolare:
			\begin{itemize}
				\item \underline{\textit{De1.1}}: Permettere di inserire o modificare commenti associati a quel documento remoto;
				\item \underline{\textit{De1.2}}: Notificare via commento, tramite l'apposito tag @, una persona condivisa con quel documento.
			\end{itemize}
		\item \underline{\textit{De2}}: Sviluppo di un servizio che permetta di generare un link per la condivisione di un documento, debitamete configurato secondo vari livelli di permessi;
		\item \underline{\textit{De3}}: Sviluppo di un servizio che generi un URL ad un \emph{Box} embed widget per permettere la visualizzazione di uno specifico documento tramite relativo visualizzatore nativo.
	\end{itemize}             % Introduzione
% !TEX encoding = UTF-8
% !TEX TS-program = pdflatex
% !TEX root = ../tesi.tex
% !TEX spellcheck = it-IT

%**************************************************************
\chapter{Il progetto di stage}
\label{cap:progetto-stage}
%**************************************************************

\intro{In questo capitolo viene descritto il progetto di stage nella sua interezza, definendone obiettivi da raggiungere e pianificazione delle attività nelle otto settimane che compongono il percorso.}

%**************************************************************
\section{Descrizione del progetto}

Descrizione del mio progetto

\subsection{Dettagli funzionalità con una tabella lunga}

\noindent Entrando più nel dettaglio le funzionalità di questa piattaforma sono le seguenti:
\begin{longtable}[c] { >{\centering\arraybackslash}p{3cm} >{\arraybackslash}p{10cm}}
		\toprule
		\centerline{\textbf{Funzione}}  & \centerline{\textbf{Descrizione}} \\
			\midrule
			\textbf{Collaborazione} & \textbf{Controllo automatico delle versioni}: TecnologiaInAnalisi inoltre salva le versioni precedenti, nel caso in cui sia necessario tornare indietro.
			\vspace*{5mm}

			\textbf{Anteprima senza download}: QuestaTecnologia supporta oltre 120 tipi di previewer file diversi, inclusi quelli in formato Word, Excel, PDF, AI, EPS, PSD e molti altri ancora, senza scaricarli. \\
			\addlinespace[0.3em]
			\midrule
			\addlinespace[0.3em]
			\textbf{Condivisione dei file} & \textbf{Basta allegati}: I file memorizzati nel cloud possono essere facilmente condivisi con altre persone, anche al di fuori del firewall dell'azienda, grazie a un semplice collegamento da inviare via e-mail o direttamente da NomeTecnologia. \\
			\addlinespace[0.3em]
			\midrule
			\addlinespace[0.3em]
			\textbf{Accesso mobile} & \textbf{Porta i tuoi file ovunque}: Condividi file in tutta sicurezza con il tuo team, i tuoi clienti o le agenzie esterne su qualsiasi dispositivo, online ed offline \\
			\addlinespace[0.3em]
			\midrule
			\addlinespace[0.3em]
			\textbf{Sicurezza} & \textbf{Controllo centralizzato delle informazioni}: Gestisci facilmente le politiche di accesso e condivisione all'interno e all'esterno della tua organizzazione
			\vspace*{5mm}

			\textbf{Protezione delle infrastrutture e delle operazioni}: Crittografia dei contenuti in transito e in archivio, con alla possibilità di creare chiavi di crittografia gestite dal cliente \\
			\addlinespace[0.3em]
			\midrule
			\addlinespace[0.3em]
			\textbf{Gestione contenuti} & \textbf{Classificazione}: Organizza i contenuti con modelli di metadati che raggruppano in maniera logica ogni file, ricorrendo ad attributi dei metadati impostati personalmente
			\vspace*{5mm}

			\textbf{Ricerca avanzata}: Trova rapidamente i file che cerchi assegnando metadati e utilizzando la ricerca avanzata per filtrare i risultati \\
			\addlinespace[0.3em]
			\midrule
			\addlinespace[0.3em]
			\textbf{Piattaforma sviluppatori} & \textbf{Dedeloper Platform}: Sviluppa app per dispositivi mobili e servizi per il Web, con la sicurezza di mantenere la conformità normativa e un livello di sicurezza enterprise e senza necessità di realizzare o gestire un livello di contenuti separato. \\
			\bottomrule\\
			\caption{\emph{TenologiaInAnalisi} - Funzionalità nel dettaglio}
			\label{tab:box_features}
\end{longtable}             % Ambiente
% !TEX encoding = UTF-8
% !TEX TS-program = pdflatex
% !TEX root = ../tesi.tex
% !TEX spellcheck = it-IT

%**************************************************************
\chapter{Analisi}
\label{cap:analisi}
%**************************************************************

\intro{Questo capitolo descrive l'analisi preliminare dell'ambiente, a monte della codifica dei servizi. Si tratta di uno studio di fattibilità del progetto in relazione a requisiti, vincoli e tecnologie da adottare.}

%**************************************************************

\section{Sotto al cofano}

Descrizione sezione

	%\begin{figure}[!h]
		%\centering
		%\includegraphics[width=\textwidth]{immagini/cap3/google.png}
		%\caption{\emph{Google} - sito ufficiale - \url{https://www.gloogle.com}}
		%\label{fig:platform}
	%\end{figure}

%**************************************************************
\newpage
\section{Autorizzazione?}

nuova sezione

	\subsection{Sotto-sezione}

	Ecco una sotto sezione!
	
%**************************************************************
             % Analisi
%\input{capitoli/capitolo-4}             % Tecno
%\input{capitoli/capitolo-5}             % Sviluppo
%\input{capitoli/capitolo-6}             % Qualifica
%\input{capitoli/capitolo-7}             % Conclusioni
% \appendix
% % !TEX encoding = UTF-8
% !TEX TS-program = pdflatex
% !TEX root = ../tesi.tex
% !TEX spellcheck = it-IT

%************************************************************** APPENDICE
\chapter{JWT}
\label{apx:jwt}
%**************************************************************

%\epigraph{Citazione}{Autore della citazione}

%**************************************************************

\section{Cos'è un jwt?}             % Appendice A

%**************************************************************
% Materiale finale
%**************************************************************
\backmatter
\printglossaries
% !TEX encoding = UTF-8
% !TEX TS-program = pdflatex
% !TEX root = ../tesi.tex
% !TEX spellcheck = it-IT

%**************************************************************
% Bibliografia
%**************************************************************

\cleardoublepage
\chapter{Bibliografia}

\paragraph{\Large{Articoli consultati}}\mbox{}

\begin{enumerate}[a.]
	\item \textbf{OAuth 2.0 Authorization Framework.}\\
	URL: \url{https://tools.ietf.org/html/rfc6749}
	\item \textbf{RTF JWT.}\\
	URL: \url{https://tools.ietf.org/html/rfc7519}
\end{enumerate}

\bigskip

\paragraph{\Large{Siti Web consultati}}\mbox{}

\begin{enumerate}[a.]
	\item \textbf{OAuth2.}\\
	URL: \url{https://oauth.net/2}
	\item \textbf{RSA cryptosystem}\\
	URL: \url{https://www.di-mgt.com.au/rsa_alg.html}
	\item \textbf{Sito ufficiale Box}\\
	URL: \url{https://www.box.com/it-it/home}
	\item \textbf{box-java-sdk}\\
	URL: \url{https://github.com/box/box-java-sdk}	
\end{enumerate}
\end{document}